% ============ 03. Entendimiento del objetivo del repositorio ============

\section{Entendimiento del objetivo del repositorio}

\subsection{Objetivo declarado}

El repositorio declara su objetivo en dos documentos con alcances distintos:

\begin{itemize}
  \item \texttt{informe\_final.md} (visión temprana): \emph{``Consolidar, validar y
        analizar los datos de investigadores reconocidos de las convocatorias 2017,
        2019 y 2021 del Minciencias, y producir un modelo dimensional y un conjunto
        de análisis descriptivos y multivariados para facilitar su uso en el
        observatorio institucional''}.
  \item \texttt{README.md} (visión final): observatorio crítico del sistema de
        reconocimiento con análisis longitudinal de \textbf{6 convocatorias
        (2013--2021)}, con foco en calidad de los datos, concentración territorial,
        brecha de género, redes de co-filiación, diversidad, productividad cruzada
        con el dataset de producción de grupos y tabla maestra de IES.
\end{itemize}

\begin{tcolorbox}[nota]
El objetivo \textbf{evolucionó durante la construcción del proyecto}: de 3
convocatorias (2017/2019/2021) se pasó a las 6 históricas (2013--2021) y se
incorporaron ejes nuevos (producción, diversidad, redes). Esta ampliación es
coherente con la naturaleza exploratoria de un observatorio, pero implica que las
métricas de las primeras fases (p.~ej. matrices de transición 2017--2019) deben
reinterpretarse a la luz del alcance ampliado.
\end{tcolorbox}

\subsection{Preguntas de investigación que el repositorio responde}

\begin{enumerate}
  \item ¿Qué tan confiable es la información que MinCiencias publica? ¿Permite
        analizar el sistema de manera comparable entre convocatorias?
  \item ¿El sistema retiene a los investigadores reconocidos entre convocatorias?
        ¿Qué porcentaje sube, baja o desaparece?
  \item ¿La capacidad investigativa está distribuida equitativamente o concentrada?
        (comparación per cápita con población DANE 2018)
  \item ¿Hay paridad de género por gran área OCDE? ¿La brecha en STEM se ha cerrado?
  \item ¿Cómo se conectan las instituciones a través de investigadores con doble
        afiliación? ¿Hay universidades-puente?
  \item ¿La composición del padrón refleja la diversidad poblacional de Colombia
        (etnia, discapacidad, víctimas del conflicto)?
  \item ¿Quién firma realmente la producción que MinCiencias mide? ¿La brecha de
        género se replica en outputs?
\end{enumerate}

\subsection{Evaluación del planteamiento del objetivo}

\begin{itemize}
  \item \textbf{Claridad}: el objetivo es comprensible y medible a través de las
        preguntas de investigación; cada pregunta mapea a un módulo de análisis
        (ver Tabla~\ref{tab:mapa_objetivo}).
  \item \textbf{Idoneidad del diseño}: el proyecto adopta el ciclo CRISP-DM y el
        patrón ``Pregunta $\rightarrow$ Hallazgo $\rightarrow$ Caveat'', alineado con
        la evaluación responsable de la ciencia (Principios de Leiden; Hicks et
        al., 2015).
  \item \textbf{Debilidad}: no existe un documento único de ``teoría de cambio'' o
        ``marco de indicadores'' que defina formalmente cada indicador, su fórmula,
        su umbral de alerta y su fuente de comparación; la formalización se
        recomienda como entregable de esta consultoría.
\end{itemize}

\begin{table}[htbp]
\centering
\caption{Mapa objetivo $\rightarrow$ módulos del repositorio}
\label{tab:mapa_objetivo}
\small
\begin{tabularx}{\linewidth}{p{6.2cm}X}
\toprule
\textbf{Pregunta de investigación} & \textbf{Módulo que la responde} \\
\midrule
Fiabilidad y comparabilidad de los datos &
\texttt{src/analisis/calidad.py}, \texttt{src/Transformacion.py}, tab 1 del dashboard \\
Retención y transición entre convocatorias &
\texttt{src/analisis/longitudinal.py}, \texttt{scripts/sprint2\_matrices\_transicion.py} \\
Concentración territorial &
\texttt{src/analisis/territorial.py} (HHI, cuotas), \texttt{sprint2\_territorial.py} \\
Brecha de género por área OCDE &
\texttt{src/analisis/genero.py}, \texttt{sprint2\_genero\_ocde.py} \\
Redes de co-filiación institucional &
\texttt{src/analisis/redes.py}, \texttt{sprint3\_redes.py}, \texttt{sprint3\_grafo\_interactivo.py} \\
Diversidad vs DANE/RUV &
\texttt{src/analisis/diversidad.py}, \texttt{sprint4\_diversidad.py} \\
Productividad cruzada &
\texttt{src/analisis/produccion.py}, \texttt{sprint5\_produccion.py}, \texttt{sprint5\_duckdb.py} \\
Ajustes de calidad, IES y geografía &
\texttt{sprint6\_*}, \texttt{src/analisis/calidad.py} \\
\bottomrule
\end{tabularx}
\end{table}

\subsection{Conclusión de la sección}

El objetivo del repositorio está \textbf{bien entendido y correctamente
descompuesto} en preguntas de investigación que los módulos efectivamente
responden. La principal recomendación es \textbf{formalizar el marco de
indicadores} (definición operativa, fórmula, fuente y límites de cada métrica)
para garantizar que la entrega final sea auditable y reproducible.
